\chapter{Related Work}\label{relatedwork}

This project sits at the intersection of three lines of research: graph neural
networks for clinical prediction from EHR data, explainability methods for GNNs,
and intrinsicly explainable models that build interpretability into the prediction itself. This chapter reviews each line and then positions the proposed model
against the closest prior work.

\section{Graph Neural Networks for EHR Prediction}

A patient record is relational by nature. A single visit ties a patient to
diagnoses, medications, laboratory results, and vital signs, and a patient
accumulates many such visits over time. This rich and heterogeneous structure has made GNNs a common
choice for clinical prediction, and a recent survey catalogues their use across
mortality prediction, readmission, and diagnosis~\cite{ossboll2024gnn}.

The reviewed work differs mainly in how the graph is built and whether time is
modelled. DynaGraph learns a dynamic graph over the sequence of visits and uses
visit-level attention to capture disease progression, reporting gains over a
range of baselines on ICU cohorts~\cite{mesinovic2026dynagraph}. GraphCare takes
a different route: it extracts a personalised knowledge graph for each patient
from large language models and external biomedical knowledge bases, then trains
a bi-attention graph network on the resulting structure~\cite{jiang2024graphcare}.
Both inject or learn structure beyond the raw record, and both report attention
weights as a form of interpretability. However, neither model extracts a concrete, minimal subgraph that drove the prediction, which is the level of explanation a clinician needs to verify a conclusion of a classifier model. 

\section{Explainability in Graph Neural Networks}

Because the mechanism of most GNNs remain ambiguous, a separate body of work develops methods to
explain them after training. A taxonomic survey groups these post-hoc methods
into instance-level and model-level approaches, with instance-level techniques further split into gradient-based, perturbation-based, decomposition-based, and surrogate families~\cite{yuan2022explainability}. GNNExplainer is a widely used instance-level method that learns an edge mask maximizing mutual information
with the prediction.

Two problems limit these methods in a clinical setting. The explanation is
produced by a mechanism separate from the model, so it is not guaranteed to
accurately reflect the reasoning that led to the prediction. Their behaviour is also
uneven: the GraphXAI benchmark introduces a synthetic generator with
ground-truth importance masks and four evaluation metrics, and shows that many
post-hoc explainers degrade on larger or denser graphs and can fail to preserve
fairness properties~\cite{agarwal2023graphxai}. This project uses the same
GraphXAI framework, but applies it to compare inherent explainability layers of the chosen models on real clinical graphs where no ground-truth mask exists.

\section{Prototype-Based Self-Explaining Models}

An alternative to post-hoc explanation is to make the model interpretable by
construction. Prototype learning follows case-based reasoning: a new input is
classified by comparing it to a set of learned representative examples. ProtoPNet
introduced this idea for image recognition with a prototype layer, and ProtGNN
extended it to graphs, using a Monte Carlo Tree Search to project each prototype
onto a real subgraph~\cite{zhang2022protgnn}. ProtGNN was evaluated on molecular
datasets such as MUTAG, so its behaviour on clinical data was left open.

ProtoEHR is the closest prior work to this project. It brings hierarchical
prototype learning to EHR data, organising prototypes at the code, visit, and
patient levels~\cite{cai2025protoehr}. Two gaps remain. ProtoEHR does not
identify which parts of a patient graph drive a prediction, and it does not
compare its built-in explanations against post-hoc methods. The model proposed
here targets both gaps by pairing a prototype layer with GraphXAI node attribution on intra-patient graphs.

\section{Positioning}

Table~\ref{tab:positioning} places the proposed models against their closest
methods. The combination that distinguishes this work is a prototype-based
classifier over intra-patient EHR graphs that also produces node-level
explanations and evaluates them against post-hoc explainers on the same data.

\begin{table}[htbp]
  \centering
  \small
  \caption{Positioning of the proposed model against the closest prior work.}
  \label{tab:positioning}
  \begin{tabular}{lllcc}
    \toprule
    Method & Interpretability & Domain & Node-level & Post-hoc \\
           & style            &        & explanation & comparison \\
    \midrule
    ProtGNN~\cite{zhang2022protgnn}        & Prototypes            & Molecular & Subgraph & No \\
    ProtoEHR~\cite{cai2025protoehr}        & Hierarchical prot.\   & EHR       & No       & No \\
    GraphCare~\cite{jiang2024graphcare}    & Bi-attention (KG)     & EHR       & No       & No \\
    DynaGraph~\cite{mesinovic2026dynagraph}& Visit attention       & EHR (ICU) & No       & No \\
    Method~C \tbd{name}                    & \tbd{}                & \tbd{}    & \tbd{}   & \tbd{} \\
    This work                              & Prototypes            & EHR (ED)  & Yes      & Yes \\
    \bottomrule
  \end{tabular}
\end{table}

The prototype family (ProtGNN, ProtoEHR) provides case-based reasoning but stops
short of pointing to the clinical entities behind a decision. The attention-based
EHR models (GraphCare, DynaGraph) expose attention weights but were not built to
compare those weights against independent post-hoc explainers. By combining a
prototype layer with GraphXAI node attribution on intra-patient graphs, this
project addresses the gap left by both groups.
